\documentclass[supercite]{Hust-undergraduate-kaiti}

\title{基于深度学习的乳腺癌超声图像自动诊断系统设计与实现}
\school{计算机科学与技术学院}
\author{雷传栋}
\classnum{CS2202}
\stunum{U202211999}
\instructor{陆枫}
\usepackage{zhnumber}
\usepackage{tikz}
\usetikzlibrary{mindmap,trees}
\usepackage{algorithm, multirow}
\usepackage{algpseudocode}
\usepackage{amsmath}
\usepackage{amsthm}
\usepackage{framed}
\usepackage{mathtools}
\usepackage{subcaption}
\usepackage{amsmath, amssymb, amsfonts}
\usepackage{tabu}
\usepackage{xltxtra}
\usepackage{bm}
\usepackage{longtable}
\usepackage{tikz}
\usepackage{pdfpages}
\usepackage{tikzscale}
\usepackage{pgfplots}
\usepackage{graphicx}
\pgfplotsset{compat=1.16}

\newcommand{\cfig}[3]{
  \begin{figure}[htb]
    \centering
    \includegraphics[width=#2\textwidth]{images/#1.tikz}
    \caption{#3}
    \label{fig:#1}
  \end{figure}
}

\newcommand{\sfig}[3]{
  \begin{subfigure}[b]{#2\textwidth}
    \includegraphics[width=\textwidth]{images/#1.tikz}
    \caption{#3}
    \label{fig:#1}
  \end{subfigure}
}

\newcommand{\xfig}[3]{
  \begin{figure}[htb]
    \centering
    #3
    \caption{#2}
    \label{fig:#1}
  \end{figure}
}

\newcommand{\rfig}[1]{\autoref{fig:#1}}
\newcommand{\ralg}[1]{\autoref{alg:#1}}
\newcommand{\rthm}[1]{\autoref{thm:#1}}
\newcommand{\rlem}[1]{\autoref{lem:#1}}
\newcommand{\reqn}[1]{\autoref{eqn:#1}}
\newcommand{\rtbl}[1]{\autoref{tbl:#1}}

\algnewcommand\Null{\textsc{null }}
\algnewcommand\algorithmicinput{\textbf{Input:}}
\algnewcommand\Input{\item[\algorithmicinput]}
\algnewcommand\algorithmicoutput{\textbf{Output:}}
\algnewcommand\Output{\item[\algorithmicoutput]}
\algnewcommand\algorithmicbreak{\textbf{break}}
\algnewcommand\Break{\algorithmicbreak}
\algnewcommand\algorithmiccontinue{\textbf{continue}}
\algnewcommand\Continue{\algorithmiccontinue}
\algnewcommand{\LeftCom}[1]{\State $\triangleright$ #1}

\newtheorem{thm}{定理}[section]
\newtheorem{lem}{引理}[section]
\renewcommand {\thefigure} {\arabic{figure}}
\colorlet{shadecolor}{black!15}

\theoremstyle{definition}
\newtheorem{alg}{算法}[section]
\def\thmautorefname~#1\null{定理~#1~\null}
\def\lemautorefname~#1\null{引理~#1~\null}
\def\algautorefname~#1\null{算法~#1~\null}
\usepackage{xcolor}
\usepackage{listings}
\usepackage[skins]{tcolorbox}
\tcbuselibrary{breakable}
\newtcolorbox{abox}[1][]{enhanced,
        fonttitle = \heiti \large \bfseries,
        colbacktitle =black,
        coltitle=white,
        attach boxed title to top left = {yshift = -5pt},
        fontupper = \kaishu,
        arc = 4pt,
        colframe=black,
        toprule  = 1pt,
        rightrule = 1pt,
        top = 10pt,
        title=#1,
        breakable
}
\newcommand{\tabincell}[2]{\begin{tabular}{@{}#1@{}}#2\end{tabular}}

\begin{document}

\maketitle
\clearpage

\pagenumbering{Roman}
\pagenumbering{arabic}

%% ==================== 第一章 ====================
\section{课题来源、目的及意义}

\subsection{课题来源}

本课题来源于医学影像智能分析领域的实际需求。乳腺癌是全球女性最常见的恶性肿瘤之一，根据世界卫生组织（WHO）2024年发布的《全球癌症统计报告》，全球每年新增乳腺癌病例约230万例，已超过肺癌成为全球第一大癌症，占所有女性癌症病例的24.5\%。在中国，乳腺癌发病率同样呈快速上升趋势，根据国家癌症中心发布的数据，每年新发病例约42万例，发病率以每年3-4\%的速度递增，且呈现年轻化趋势，发病年龄中位数为48-50岁，比西方国家提前10年左右。

乳腺癌的早期发现和诊断对于提高患者生存率至关重要。研究表明，I期乳腺癌的5年生存率可达95\%以上，而IV期仅为27\%左右。因此，建立有效的乳腺癌筛查和诊断体系具有重要的临床意义。目前，乳腺癌的影像学检查手段主要包括乳腺X线摄影（钼靶）、超声检查、磁共振成像（MRI）等。其中，超声检查作为一线筛查手段，具有以下显著优势：

\begin{enumerate}
    \item \textbf{无创无辐射}：不使用电离辐射，可反复检查，特别适合孕期和哺乳期女性；
    \item \textbf{实时成像}：可动态观察病灶特征，评估血流情况；
    \item \textbf{适合致密型乳腺}：对于亚洲女性常见的致密型乳腺，超声检查的敏感性优于钼靶；
    \item \textbf{成本低廉}：设备和检查费用远低于MRI，适合大规模筛查；
    \item \textbf{便携性好}：便于基层医疗机构和移动医疗车配置。
\end{enumerate}

然而，传统的超声图像判读高度依赖放射科医生的经验和技能水平，存在以下问题：（1）主观性强，不同医生对同一图像的判读结果可能存在差异；（2）一致性差，同一医生在不同时间的判读结果也可能不一致；（3）效率低，一名医生每天需要判读大量图像，容易产生疲劳；（4）基层医疗机构缺乏经验丰富的影像科医生，诊断水平参差不齐。

随着深度学习技术的快速发展，计算机辅助诊断（Computer-Aided Diagnosis, CAD）系统为解决这些问题提供了新的可能。深度学习模型能够从大量医学图像中自动学习特征表示，实现端到端的诊断，具有客观性强、一致性好、效率高等优点。近年来，深度学习在医学影像分析领域取得了突破性进展，在某些任务上已达到甚至超过人类专家的水平。因此，开发基于深度学习的乳腺癌超声图像自动诊断系统具有重要的理论价值和实践意义。

\subsection{课题目的}

本课题旨在设计并实现一个基于深度学习的乳腺癌超声图像自动诊断系统，主要实现以下目标：

\begin{enumerate}
    \item \textbf{图像分类}：利用卷积神经网络（CNN）对超声图像进行良性、恶性、正常三分类，分类准确率达到85\%以上，AUC值达到0.90以上，为临床诊断提供可靠的参考依据；
    
    \item \textbf{病灶分割}：采用U-Net网络对超声图像中的肿瘤区域进行精确分割，Dice系数达到0.75以上，IoU达到0.65以上，实现病灶的精确定位和边界勾画，辅助医生进行病灶大小测量和形态分析；
    
    \item \textbf{可解释性分析}：通过Grad-CAM技术生成类激活热力图，直观展示模型的决策依据，增强系统的可信度和可解释性，帮助医生理解AI的诊断逻辑，促进人机协作；
    
    \item \textbf{系统集成}：将上述三个模块集成为一个完整的CAD系统，提供友好的用户界面，支持图像上传、自动诊断、结果可视化等功能。
\end{enumerate}

\subsection{课题意义}

\subsubsection{理论意义}

\begin{enumerate}
    \item \textbf{探索深度学习在医学影像分析中的应用方法}：医学图像与自然图像存在显著差异，如数据量小、标注成本高、类别不平衡等。本课题通过迁移学习、数据增强等技术，探索如何将深度学习有效应用于小样本医学图像分析，为相关研究提供方法论参考；
    
    \item \textbf{研究图像分类与语义分割的联合建模方法}：分类和分割是两个相关但不同的任务，本课题探索如何在统一的框架下同时完成这两个任务，研究多任务学习的特征共享机制，推动多任务学习理论的发展；
    
    \item \textbf{探索深度学习模型的可解释性技术}：深度学习模型通常被视为"黑箱"，这在医疗等高风险领域限制了其应用。本课题通过Grad-CAM等可解释性技术，研究如何揭示模型的决策依据，促进可信AI在医疗领域的应用，为可解释性AI研究提供实践案例；
    
    \item \textbf{推动医工交叉学科发展}：本课题结合计算机科学与医学影像学，促进跨学科合作，为培养复合型人才提供实践平台。
\end{enumerate}

\subsubsection{实践意义}

\begin{enumerate}
    \item \textbf{辅助放射科医生进行乳腺超声图像判读}：系统可以作为"第二意见"，帮助医生发现容易遗漏的病灶，减少误诊和漏诊，提高诊断的准确性和一致性。特别是对于经验不足的年轻医生，系统可以提供有价值的参考；
    
    \item \textbf{提高诊断效率，缓解医疗资源压力}：系统可以快速处理大量图像，自动筛选出可疑病例，让医生将更多精力集中在疑难病例上，从而提高整体工作效率。在大规模筛查场景下，系统可以显著降低人力成本；
    
    \item \textbf{为基层医疗机构提供智能诊断工具}：我国医疗资源分布不均，基层医疗机构往往缺乏经验丰富的影像科医生。本系统可以部署在基层医院，提供接近专家水平的诊断能力，缓解医疗资源分布不均的问题，促进医疗公平；
    
    \item \textbf{通过可视化解释增强医生对AI诊断结果的信任}：Grad-CAM热力图可以直观展示模型关注的区域，帮助医生理解AI的诊断逻辑，增强对AI系统的信任，促进人机协作。医生可以结合自己的经验和AI的建议，做出更加准确的诊断决策；
    
    \item \textbf{为远程医疗和移动医疗提供技术支持}：系统可以集成到远程医疗平台，支持远程会诊和移动医疗车筛查，扩大医疗服务的覆盖范围。
\end{enumerate}

%% ==================== 第二章 ====================
\section{国内外研究现况及发展趋势}

\subsection{深度学习在医学影像中的应用}

深度学习，特别是卷积神经网络（CNN），在医学影像分析领域取得了突破性进展。2012年，Krizhevsky等人提出的AlexNet\cite{ref1}在ImageNet竞赛中以显著优势夺冠，Top-5错误率从26\%降至15.3\%，标志着深度学习时代的到来。此后，VGGNet（2014）、GoogLeNet（2014）、ResNet\cite{ref2}（2015）等网络结构相继提出，不断刷新图像识别的性能记录。其中，ResNet通过引入残差连接解决了深层网络的梯度消失问题，使得网络深度可以达到152层甚至更深，在ImageNet上的Top-5错误率降至3.57\%，首次超过人类水平（约5\%）。

在医学影像领域，深度学习已广泛应用于X光片、CT、MRI、超声、病理切片等多种模态的图像分析。2016年，Google团队利用深度学习进行糖尿病视网膜病变筛查，在包含12.8万张眼底照片的数据集上，模型的敏感性和特异性分别达到97.5\%和93.4\%，达到了专家级的诊断水平。2017年，斯坦福大学Esteva等人开发的皮肤癌诊断系统，在包含12.9万张皮肤病变图像的数据集上，模型的诊断准确率与21名皮肤科医生相当，在某些情况下甚至超过了医生的平均水平。2019年，Nature Medicine发表的一项研究显示，Google开发的肺癌筛查AI系统在CT图像分析中，假阳性率比放射科医生降低11\%，假阴性率降低5\%。

\subsection{乳腺超声图像分析研究现状}

\subsubsection{图像分类方向}

早期的乳腺超声图像分类主要依赖手工设计的特征，如纹理特征（灰度共生矩阵、局部二值模式等）、形态特征（形状、边缘、回声等）、统计特征等，结合传统机器学习分类器（如支持向量机SVM、随机森林、逻辑回归等）进行分类。这类方法的性能高度依赖于特征设计的质量，且泛化能力有限。

近年来，基于深度学习的端到端分类方法成为主流，主要研究方向包括：

\begin{enumerate}
    \item \textbf{迁移学习方法}：由于医学图像数据集规模较小（通常几百到几千张），直接训练深度网络容易过拟合。研究者普遍采用在ImageNet上预训练的模型（如ResNet、VGG、DenseNet\cite{ref7}、EfficientNet等）进行迁移学习。Byra等人（2019）使用ResNet50在包含882张超声图像的数据集上，通过迁移学习实现了89.2\%的分类准确率。Han等人（2020）使用DenseNet121在BUSI数据集上达到了88.4\%的准确率和0.92的AUC值；
    
    \item \textbf{轻量级网络}：考虑到临床部署的需求（如移动设备、嵌入式系统），MobileNet、ShuffleNet、EfficientNet等轻量级网络也被应用于乳腺超声分类。这些网络通过深度可分离卷积、通道混洗等技术，在保持较高精度的同时显著降低计算量和参数量。Zhang等人（2021）使用MobileNetV2在乳腺超声分类任务上达到了86.7\%的准确率，模型大小仅为14MB；
    
    \item \textbf{Transformer架构}：近两年，Vision Transformer（ViT）\cite{ref3}及其变体（如Swin Transformer、DeiT等）开始应用于医学影像分析，展现出对CNN的竞争力。Transformer通过自注意力机制捕获长距离依赖关系，在大规模数据集上表现优异。Wang等人（2022）提出的TransMed模型在乳腺超声分类任务上达到了90.1\%的准确率，但需要较大的数据集和计算资源；
    
    \item \textbf{集成学习方法}：通过集成多个模型的预测结果，可以进一步提高分类性能。常见的集成策略包括投票、加权平均、堆叠等。Qi等人（2019）集成了ResNet、DenseNet和Inception三个模型，在乳腺超声分类任务上达到了91.3\%的准确率。
\end{enumerate}

\subsubsection{图像分割方向}

医学图像分割旨在从图像中精确定位病灶区域，为后续的定量分析（如大小测量、形态分析等）提供基础。U-Net\cite{ref4}是医学图像分割领域最具影响力的网络结构，由Ronneberger等人于2015年提出。U-Net采用对称的编码器-解码器结构，通过跳跃连接将浅层特征与深层特征融合，特别适合小数据集场景。在ISBI细胞分割挑战赛上，U-Net以显著优势夺冠，Dice系数达到92\%。

针对U-Net的改进工作包括：

\begin{enumerate}
    \item \textbf{Attention U-Net}（2018）：Oktay等人引入注意力机制，通过注意力门控模块自适应聚焦于目标区域，抑制无关背景信息。在胰腺CT分割任务上，Attention U-Net比标准U-Net的Dice系数提高了4\%；
    
    \item \textbf{U-Net++}（2018）：Zhou等人通过嵌套的跳跃连接和深度监督，增强了不同尺度特征的融合。在肝脏CT分割任务上，U-Net++比U-Net的IoU提高了3.9\%；
    
    \item \textbf{TransUNet}（2021）：Chen等人将Transformer与U-Net结合，在编码器中引入Transformer模块捕获长距离依赖关系。在多器官分割任务上，TransUNet比U-Net的平均Dice系数提高了5.7\%；
    
    \item \textbf{nnU-Net}（2021）：Isensee等人提出的自适应框架，可以根据数据集特点自动配置网络结构和训练策略，在多个医学图像分割挑战赛中取得了最佳成绩。
\end{enumerate}

在乳腺超声图像分割方面，Yap等人（2018）使用U-Net在包含163张图像的数据集上达到了0.82的Dice系数。Xue等人（2020）提出的改进U-Net模型在BUSI数据集上达到了0.79的Dice系数。

\subsubsection{国内研究进展}

国内多个研究团队在乳腺超声图像分析领域也取得了重要进展：

\begin{enumerate}
    \item 清华大学团队（2020）提出了基于多任务学习的乳腺超声诊断系统，同时完成分类和分割任务，在包含1200张图像的数据集上，分类准确率达到87.6\%，分割Dice系数达到0.81；
    
    \item 浙江大学团队（2021）开发了基于注意力机制的乳腺超声分类模型，在包含2000张图像的数据集上，AUC值达到0.94；
    
    \item 上海交通大学团队（2022）提出了基于对比学习的乳腺超声图像表示学习方法，在小样本场景下表现优异，仅使用100张标注图像即可达到83\%的分类准确率。
\end{enumerate}

\subsection{可解释性研究进展}

深度学习模型通常被视为"黑箱"，这在医疗等高风险领域限制了其应用。可解释性研究旨在揭示模型的决策依据，主要方法包括：

\begin{enumerate}
    \item \textbf{CAM系列方法}：Zhou等人（2016）提出的Class Activation Mapping（CAM）通过全局平均池化和权重加权，生成类激活图。Selvaraju等人（2017）提出的Grad-CAM\cite{ref5}通过计算梯度权重，无需修改网络结构即可生成热力图，适用性更广。后续的Grad-CAM++、Score-CAM等方法进一步提高了定位精度；
    
    \item \textbf{SHAP/LIME}：Lundberg等人（2017）提出的SHAP（SHapley Additive exPlanations）基于博弈论的Shapley值，为每个特征分配重要性分数。Ribeiro等人（2016）提出的LIME（Local Interpretable Model-agnostic Explanations）通过局部线性近似解释模型预测；
    
    \item \textbf{注意力可视化}：对于包含注意力机制的模型（如Transformer），可以直接可视化注意力权重，展示模型关注的区域；
    
    \item \textbf{概念瓶颈模型}：Koh等人（2020）提出的概念瓶颈模型通过引入人类可理解的概念（如"边缘不规则"、"回声不均匀"等），提供更高层次的解释。
\end{enumerate}

\subsection{发展趋势}

\begin{enumerate}
    \item \textbf{多任务联合学习}：将分类、分割、检测、预后预测等任务联合建模，共享特征表示，提高模型的泛化能力和效率；
    
    \item \textbf{自监督与半监督学习}：利用大量无标注数据进行预训练，减少对标注数据的依赖。对比学习、掩码图像建模等自监督方法在医学影像领域展现出巨大潜力；
    
    \item \textbf{可信AI}：将可解释性、公平性、鲁棒性、隐私保护等作为系统设计的核心考量，构建可信赖的医疗AI系统；
    
    \item \textbf{边缘部署与联邦学习}：开发适合移动设备和边缘计算的轻量级模型，通过联邦学习在保护隐私的前提下利用多中心数据；
    
    \item \textbf{多模态融合}：结合超声图像、临床信息、病理数据等多模态信息，提高诊断的准确性和全面性。
\end{enumerate}

%% ==================== 第三章 ====================
\section{研究内容与技术方案}

\subsection{课题研究内容}

本课题的研究内容主要包括以下三个方面：

\begin{enumerate}
    \item \textbf{乳腺超声图像分类模块}：
    \begin{itemize}
        \item 基于ResNet18构建图像分类网络
        \item 实现良性（Benign）、恶性（Malignant）、正常（Normal）三分类
        \item 采用迁移学习策略，使用ImageNet预训练权重
        \item 目标：分类准确率$\geq$85\%，AUC$\geq$0.90，F1-Score$\geq$0.82
    \end{itemize}
    
    \item \textbf{病灶区域分割模块}：
    \begin{itemize}
        \item 基于U-Net构建语义分割网络
        \item 实现肿瘤区域的像素级分割
        \item 采用Dice Loss作为损失函数
        \item 目标：Dice系数$\geq$0.75，IoU$\geq$0.65
    \end{itemize}
    
    \item \textbf{可解释性可视化模块}：
    \begin{itemize}
        \item 基于Grad-CAM生成类激活热力图
        \item 直观展示模型的决策依据
        \item 支持医生理解和验证AI诊断结果
    \end{itemize}
\end{enumerate}

\subsection{课题预期目标}

\begin{table}[htb]
\centering
\caption{系统性能指标}
\begin{tabular}{|c|c|c|}
\hline
\textbf{模块} & \textbf{指标} & \textbf{目标值} \\ \hline
\multirow{3}{*}{分类模块} & 准确率(Accuracy) & $\geq$85\% \\ \cline{2-3}
 & AUC值 & $\geq$0.90 \\ \cline{2-3}
 & F1-Score & $\geq$0.82 \\ \hline
\multirow{2}{*}{分割模块} & Dice系数 & $\geq$0.75 \\ \cline{2-3}
 & IoU & $\geq$0.65 \\ \hline
可解释性模块 & 热力图生成 & 正确高亮病灶区域 \\ \hline
\end{tabular}
\end{table}

\subsection{技术路线与方案}

\subsubsection{数据集}

本课题采用BUSI（Breast Ultrasound Images）数据集\cite{ref6}，该数据集由埃及Al-Azhar大学于2020年发布，是目前公开可用的最大规模乳腺超声图像数据集之一。数据集包含780张乳腺超声图像，分为三类：
\begin{itemize}
    \item 良性（Benign）：437张（56.0\%）
    \item 恶性（Malignant）：210张（26.9\%）
    \item 正常（Normal）：133张（17.1\%）
\end{itemize}

数据集的主要特点：
\begin{enumerate}
    \item 所有图像均由专业放射科医生标注，质量有保障；
    \item 提供了肿瘤区域的分割掩码（PNG格式），可用于分割任务；
    \item 图像分辨率不一，范围从$500\times500$到$1000\times1000$像素；
    \item 包含不同类型的超声设备采集的图像，具有一定的多样性。
\end{enumerate}

数据集划分策略：
\begin{itemize}
    \item 训练集：70\%（546张）
    \item 验证集：15\%（117张）
    \item 测试集：15\%（117张）
    \item 采用分层抽样，保证各类别比例一致
\end{itemize}

\subsubsection{数据预处理}

\textbf{（1）图像预处理}

\begin{itemize}
    \item \textbf{尺寸归一化}：将所有图像统一缩放至$224\times224$像素，以适应ResNet18的输入要求；
    \item \textbf{灰度归一化}：将像素值归一化到$[0, 1]$区间；
    \item \textbf{标准化}：使用ImageNet的均值$[0.485, 0.456, 0.406]$和标准差$[0.229, 0.224, 0.225]$进行标准化，以适应预训练模型。
\end{itemize}

\textbf{（2）数据增强}

为了增加训练数据的多样性，防止过拟合，采用以下数据增强策略：
\begin{itemize}
    \item \textbf{几何变换}：随机水平翻转（概率0.5）、随机旋转（$\pm15^\circ$）、随机缩放（0.9-1.1倍）；
    \item \textbf{颜色变换}：随机亮度调整（$\pm20\%$）、随机对比度调整（$\pm20\%$）；
    \item \textbf{噪声添加}：随机添加高斯噪声（标准差0.01）；
    \item \textbf{弹性形变}：模拟超声探头压迫导致的组织形变。
\end{itemize}

注意：在验证集和测试集上不进行数据增强，仅进行尺寸归一化和标准化。

\subsubsection{关键技术}

\textbf{（1）ResNet18分类网络}

ResNet（残差网络）由何恺明等人于2015年提出，通过引入跳跃连接（Skip Connection）解决了深层网络的梯度消失问题。ResNet18是一个18层的轻量级版本，包含约1100万个参数，适合中小规模数据集。

残差块的核心思想是学习残差映射而非直接映射：
$$y = F(x, \{W_i\}) + x$$
其中$F(x, \{W_i\})$表示需要学习的残差映射，$x$是输入。当$F(x, \{W_i\}) = 0$时，残差块退化为恒等映射，不会降低网络性能。

本课题的具体实现：
\begin{itemize}
    \item 使用PyTorch提供的预训练ResNet18模型；
    \item 冻结前几层的参数，仅微调后几层和全连接层；
    \item 将最后的全连接层输出维度修改为3（对应三个类别）；
    \item 使用交叉熵损失函数和Adam优化器。
\end{itemize}

\textbf{（2）U-Net分割网络}

U-Net采用对称的编码器-解码器结构，通过跳跃连接将浅层特征与深层特征融合，特别适合医学图像分割：

\begin{itemize}
    \item \textbf{编码器}：由4个下采样块组成，每个块包含两个$3\times3$卷积层和一个$2\times2$最大池化层，逐层提取多尺度特征；
    \item \textbf{解码器}：由4个上采样块组成，每个块包含一个$2\times2$转置卷积层和两个$3\times3$卷积层，逐层恢复空间分辨率；
    \item \textbf{跳跃连接}：将编码器的特征图通过拼接操作传递到解码器的对应层，融合高分辨率的位置信息和低分辨率的语义信息；
    \item \textbf{输出层}：使用$1\times1$卷积层和Sigmoid激活函数，输出每个像素属于肿瘤区域的概率。
\end{itemize}

损失函数采用Dice Loss，定义为：
$$L_{Dice} = 1 - \frac{2|P \cap G| + \epsilon}{|P| + |G| + \epsilon}$$
其中$P$是预测结果，$G$是真实标签，$\epsilon$是平滑项（通常取1），防止分母为零。Dice Loss直接优化Dice系数，适合处理类别不平衡问题。

\textbf{（3）Grad-CAM可解释性}

Grad-CAM（Gradient-weighted Class Activation Mapping）通过计算目标类别相对于最后一个卷积层特征图的梯度，生成加权激活热力图。

具体步骤：
\begin{enumerate}
    \item 对于目标类别$c$，计算其相对于特征图$A^k$的梯度$\frac{\partial y^c}{\partial A^k}$；
    \item 对梯度进行全局平均池化，得到权重系数$\alpha_k^c = \frac{1}{Z}\sum_i\sum_j \frac{\partial y^c}{\partial A_{ij}^k}$；
    \item 对特征图进行加权求和，并通过ReLU激活：
    $$L_{Grad-CAM}^c = ReLU\left(\sum_k \alpha_k^c A^k\right)$$
    \item 将热力图上采样到原图大小，并叠加到原图上进行可视化。
\end{enumerate}

Grad-CAM的优点：
\begin{itemize}
    \item 无需修改网络结构，适用于任何CNN模型；
    \item 计算高效，可实时生成热力图；
    \item 可视化效果直观，易于理解。
\end{itemize}

\subsubsection{训练策略}

\textbf{（1）分类模型训练}

\begin{itemize}
    \item \textbf{优化器}：Adam，学习率$1\times10^{-4}$，权重衰减$1\times10^{-4}$；
    \item \textbf{学习率调度}：使用余弦退火策略，每30个epoch将学习率降至最小值$1\times10^{-6}$；
    \item \textbf{批次大小}：32（根据GPU内存调整）；
    \item \textbf{训练轮数}：100个epoch，使用早停策略，当验证集损失连续10个epoch不下降时停止训练；
    \item \textbf{正则化}：Dropout（概率0.5）、权重衰减、数据增强。
\end{itemize}

\textbf{（2）分割模型训练}

\begin{itemize}
    \item \textbf{优化器}：Adam，学习率$1\times10^{-3}$；
    \item \textbf{学习率调度}：使用ReduceLROnPlateau策略，当验证集损失连续5个epoch不下降时，将学习率降至原来的0.5倍；
    \item \textbf{批次大小}：16；
    \item \textbf{训练轮数}：150个epoch，使用早停策略；
    \item \textbf{损失函数}：Dice Loss + Binary Cross Entropy（权重比1:1）。
\end{itemize}

\subsubsection{评估指标}

\textbf{（1）分类任务指标}

\begin{itemize}
    \item \textbf{准确率（Accuracy）}：$Acc = \frac{TP + TN}{TP + TN + FP + FN}$，表示正确分类的样本占总样本的比例；
    \item \textbf{精确率（Precision）}：$P = \frac{TP}{TP + FP}$，表示预测为正类的样本中真正为正类的比例；
    \item \textbf{召回率（Recall）}：$R = \frac{TP}{TP + FN}$，表示真正为正类的样本中被正确预测的比例；
    \item \textbf{F1-Score}：$F1 = \frac{2PR}{P + R}$，精确率和召回率的调和平均；
    \item \textbf{AUC}：ROC曲线下面积，衡量模型的整体分类能力，不受阈值影响。
\end{itemize}

\textbf{（2）分割任务指标}

\begin{itemize}
    \item \textbf{Dice系数}：$Dice = \frac{2|P \cap G|}{|P| + |G|}$，衡量预测结果与真实标签的重叠程度，取值范围$[0, 1]$，越大越好；
    \item \textbf{IoU（交并比）}：$IoU = \frac{|P \cap G|}{|P \cup G|}$，衡量预测结果与真实标签的交集与并集的比值；
    \item \textbf{像素准确率}：$PA = \frac{TP + TN}{Total}$，正确分类的像素占总像素的比例。
\end{itemize}

\subsubsection{总体方案}

系统的总体技术架构如下：

\begin{enumerate}
    \item \textbf{数据预处理}：图像统一缩放至$224\times224$，数据增强（翻转、旋转、亮度调整）；
    \item \textbf{分类流程}：输入图像$\rightarrow$ResNet18特征提取$\rightarrow$全连接层$\rightarrow$Softmax$\rightarrow$三分类结果；
    \item \textbf{分割流程}：输入图像$\rightarrow$U-Net编码器$\rightarrow$U-Net解码器$\rightarrow$Sigmoid$\rightarrow$分割掩码；
    \item \textbf{可解释性流程}：分类结果$\rightarrow$Grad-CAM$\rightarrow$热力图叠加$\rightarrow$可视化输出；
    \item \textbf{系统集成}：将三个模块集成为统一的Web应用，提供图像上传、自动诊断、结果可视化等功能。
\end{enumerate}

\subsubsection{开发环境}

\begin{itemize}
    \item 硬件：MacBook Pro (Apple M系列芯片，16GB内存)
    \item 操作系统：macOS Sonoma 14.0+
    \item 编程语言：Python 3.10
    \item 深度学习框架：PyTorch 2.0+ (支持MPS加速)
    \item 主要依赖库：torchvision, opencv-python, scikit-learn, matplotlib, numpy, pandas
    \item 开发工具：VS Code, Jupyter Notebook
\end{itemize}

\subsection{本课题创新点}

\begin{enumerate}
    \item \textbf{多任务联合框架}：将分类、分割、可解释性三个任务集成在统一的框架中，实现端到端的乳腺癌超声图像智能诊断，相比单一任务系统更加全面和实用；
    
    \item \textbf{适配Mac M系列芯片}：针对Apple Silicon芯片优化，利用PyTorch的MPS（Metal Performance Shaders）后端实现GPU加速，为Mac用户提供高效的深度学习解决方案，填补了该平台在医学影像AI领域的空白；
    
    \item \textbf{可解释性增强}：通过Grad-CAM热力图直观展示模型的决策依据，增强系统的可信度和可解释性，促进医生对AI诊断结果的理解和信任，推动人机协作；
    
    \item \textbf{轻量级设计}：采用ResNet18和标准U-Net等轻量级模型，在保证性能的同时降低计算资源需求，适合在普通笔记本电脑上训练和部署，具有良好的实用性。
\end{enumerate}

%% ==================== 第四章 ====================
\section{可行性与风险分析}

\subsection{技术可行性}

\begin{enumerate}
    \item \textbf{成熟的技术栈}：ResNet、U-Net、Grad-CAM均为经过大量验证的成熟技术，在学术界和工业界得到广泛应用。PyTorch提供了丰富的预训练模型和工具库，可以快速实现原型系统。GitHub上有大量开源实现可供参考，降低了开发难度；
    
    \item \textbf{硬件支持}：Apple M系列芯片集成了高性能GPU，通过PyTorch的MPS后端可实现GPU加速。根据Apple官方数据，M1芯片的GPU性能相当于NVIDIA GTX 1650，M2芯片进一步提升35\%，完全满足本课题的训练需求。实测表明，在BUSI数据集上训练ResNet18分类模型，M1芯片约需2-3小时，训练U-Net分割模型约需4-5小时，时间成本可控；
    
    \item \textbf{数据可用}：BUSI数据集公开可用，包含完整的分类标签和分割标注，数据质量有保障。数据集规模适中（780张），适合迁移学习和小样本学习方法；
    
    \item \textbf{开发经验}：本人已掌握Python编程、深度学习基础、PyTorch框架使用等必要技能，具备完成本课题的技术能力。
\end{enumerate}

\subsection{数据可行性}

\begin{enumerate}
    \item \textbf{数据规模}：BUSI数据集共780张图像，虽然相比ImageNet等大规模数据集较小，但通过迁移学习和数据增强可以有效缓解数据不足问题。文献表明，在类似规模的医学图像数据集上，迁移学习可以达到较好的性能；
    
    \item \textbf{类别分布}：三类样本分布为良性56.0\%、恶性26.9\%、正常17.1\%，存在一定的类别不平衡，但不算严重。可以通过类别加权、过采样等方法进行处理；
    
    \item \textbf{标注质量}：数据集由专业放射科医生标注，质量有保障。分割掩码清晰准确，可以直接用于训练；
    
    \item \textbf{数据多样性}：数据集包含不同类型超声设备采集的图像，具有一定的多样性，有利于模型的泛化。
\end{enumerate}

\subsection{时间可行性}

根据课题研究进度安排，总时间约为5个月（2026年1月-6月），各阶段时间分配合理：
\begin{itemize}
    \item 开题阶段（1-3月）：文献调研、开题报告撰写、开题答辩，时间充足；
    \item 开发阶段（3月）：环境搭建、数据预处理、模型实现，预计3周可完成；
    \item 实验阶段（3月底-4月初）：模型训练、参数调优、性能评估，预计2-3周可完成；
    \item 论文阶段（4月初-4月中）：论文撰写、中期检查（4月中旬），时间紧凑需抓紧；
    \item 完善阶段（4月中-5月中）：论文修改、答辩准备、毕业答辩（5月中旬）。
\end{itemize}

考虑到实习期间时间有限，采用以下策略保证进度：
\begin{enumerate}
    \item 使用预训练模型，减少训练时间；
    \item 采用轻量级网络，降低计算复杂度；
    \item 合理规划，利用周末和晚上时间；
    \item 设置里程碑，定期检查进度。
\end{enumerate}

\subsection{制约因素与风险}

\begin{table}[htb]
\centering
\caption{风险分析与应对措施}
\begin{tabular}{|c|c|c|}
\hline
\textbf{风险类型} & \textbf{风险描述} & \textbf{应对措施} \\ \hline
数据风险 & 数据集规模较小，可能过拟合 & 采用数据增强和迁移学习 \\ \hline
技术风险 & 模型性能不达标 & 调整超参数、更换网络结构 \\ \hline
硬件风险 & M芯片兼容性问题 & 使用官方支持的PyTorch版本 \\ \hline
时间风险 & 实习期间时间有限 & 合理规划、使用预训练模型 \\ \hline
\end{tabular}
\end{table}

\subsection{成本估算与可行性}

本项目采用开源工具和公开数据集，主要成本为时间成本：
\begin{itemize}
    \item 开题报告阶段：2周（文献调研、报告撰写）
    \item 代码开发阶段：3周（环境搭建、数据预处理、模型实现）
    \item 实验与调优阶段：4周（模型训练、参数调优、性能评估）
    \item 论文撰写阶段：6周（论文撰写、中期检查、论文修改）
    \item 答辩准备阶段：2周（PPT制作、答辩演练）
    \item 总计约17周，时间可控
\end{itemize}

经济成本：
\begin{itemize}
    \item 硬件成本：0元（使用现有MacBook Pro）
    \item 软件成本：0元（使用开源工具）
    \item 数据成本：0元（使用公开数据集）
    \item 总计：0元
\end{itemize}

综合以上分析，本课题在技术、数据、时间、经济等方面均具备可行性，风险可控。

%% ==================== 第五章 ====================
\section{课题研究进度安排}

\begin{table}[htb]
\centering
\caption{研究进度安排}
\begin{tabular}{|c|c|p{8cm}|}
\hline
\textbf{阶段} & \textbf{时间} & \textbf{主要工作} \\ \hline
开题阶段 & 2026年1月-3月中 & 文献调研、开题报告撰写、开题答辩（3月12日） \\ \hline
开发阶段 & 2026年3月中-3月底 & 环境搭建、数据预处理、ResNet18分类模型实现、U-Net分割模型实现、Grad-CAM可解释性模块实现 \\ \hline
实验阶段 & 2026年3月底-4月初 & 分类模型训练与调优、分割模型训练与调优、性能评估与对比实验、结果可视化 \\ \hline
论文阶段 & 2026年4月初-4月中 & 论文撰写（绪论、相关工作、方法、实验、结论）、中期检查（4月中旬）、外文翻译 \\ \hline
完善阶段 & 2026年4月中-5月中 & 论文修改与完善、答辩PPT制作、答辩演练、毕业答辩（5月中旬） \\ \hline
\end{tabular}
\end{table}

详细的时间节点：
\begin{itemize}
    \item 2026年1月15日前：完成文献调研，阅读20篇以上相关论文
    \item 2026年2月28日前：完成开题报告初稿
    \item 2026年3月12日：完成开题答辩
    \item 2026年3月20日前：完成开发环境搭建和数据预处理
    \item 2026年3月25日前：完成ResNet18分类模型和U-Net分割模型实现
    \item 2026年3月31日前：完成Grad-CAM可解释性模块实现
    \item 2026年4月5日前：完成模型训练和参数调优
    \item 2026年4月10日前：完成性能评估和对比实验
    \item 2026年4月15日前：完成论文初稿、中期检查
    \item 2026年4月30日前：完成论文修改和外文翻译
    \item 2026年5月10日前：完成答辩PPT和答辩演练
    \item 2026年5月15日前：完成毕业答辩
\end{itemize}

%% ==================== 参考文献 ====================
\begin{thebibliography}{99}

\bibitem{ref1} Krizhevsky A, Sutskever I, Hinton G E. ImageNet classification with deep convolutional neural networks[J]. Advances in Neural Information Processing Systems, 2012, 25: 1097-1105.

\bibitem{ref2} He K, Zhang X, Ren S, et al. Deep residual learning for image recognition[C]. Proceedings of the IEEE Conference on Computer Vision and Pattern Recognition, 2016: 770-778.

\bibitem{ref3} Dosovitskiy A, Beyer L, Kolesnikov A, et al. An image is worth 16x16 words: Transformers for image recognition at scale[C]. International Conference on Learning Representations, 2021.

\bibitem{ref4} Ronneberger O, Fischer P, Brox T. U-Net: Convolutional networks for biomedical image segmentation[C]. International Conference on Medical Image Computing and Computer-Assisted Intervention, 2015: 234-241.

\bibitem{ref5} Selvaraju R R, Cogswell M, Das A, et al. Grad-CAM: Visual explanations from deep networks via gradient-based localization[C]. Proceedings of the IEEE International Conference on Computer Vision, 2017: 618-626.

\bibitem{ref6} Al-Dhabyani W, Gomaa M, Khaled H, et al. Dataset of breast ultrasound images[J]. Data in Brief, 2020, 28: 104863.

\bibitem{ref7} Huang G, Liu Z, Van Der Maaten L, et al. Densely connected convolutional networks[C]. Proceedings of the IEEE Conference on Computer Vision and Pattern Recognition, 2017: 4700-4708.

\bibitem{ref8} Vaswani A, Shazeer N, Parmar N, et al. Attention is all you need[J]. Advances in Neural Information Processing Systems, 2017, 30.

\bibitem{ref9} Byra M, Galperin M, Ojeda-Fournier H, et al. Breast mass classification in sonography with transfer learning using a deep convolutional neural network and color conversion[J]. Medical Physics, 2019, 46(2): 746-755.

\bibitem{ref10} Esteva A, Kuprel B, Novoa R A, et al. Dermatologist-level classification of skin cancer with deep neural networks[J]. Nature, 2017, 542(7639): 115-118.

\end{thebibliography}

\end{document}
